% Lean module: MonadFormal.Sandbox
\subsubsection{Present and continuation-sensitive distinctions}
\label{ss:continuations}

For a context $C$, present indistinguishability is
\begin{equation}
  x \presentEq{C} y
  \iff p(x)=p(y)\quad\text{for every }p\in\Probes{C}.
  \label{eq:present-eq}
\end{equation}
This is equality of admitted observations, not equality of entities.
Let $\Conts{C}$ be an admitted family of continuations $k:C\to D$.
Their continuation-probe family is
\begin{equation}
  \ContProbes{C}=\{p\circ k_*:k\in\Conts{C},\;
    p\in\Probes{\operatorname{cod}(k)}\}.
  \label{eq:cont-probe-family}
\end{equation}
Define
\begin{equation}
  x\contEq{C}y\iff f(x)=f(y)
    \quad\text{for every }f\in\ContProbes{C}.
  \label{eq:cont-eq}
\end{equation}
The future enters through admitted tests of possible continuations;
no information from an actual future is assumed available.

Both relations use the same mathematical object: an indexed family
$F=(f_i:X\to O_i)_{i\in I}$, with possibly different codomains.
Write $x\approx_F y$ when $f_i(x)=f_i(y)$ for every $i$.
\leandef{MonadFormal.ProbeFamily,MonadFormal.indist}

\begin{theorem}[Observational equivalence]\label{thm:observational-equivalence}
For every probe family $F$, the relation $\approx_F$ is an equivalence
relation on $X$.
\end{theorem}
\begin{proof}
For each $i$, equality in $O_i$ is reflexive, symmetric and transitive.
Applying the corresponding equality law separately at every index proves
the three laws for $\approx_F$.
\leanproof{MonadFormal.indist_refl,MonadFormal.indist_symm,MonadFormal.indist_trans}
\end{proof}
